\section{Rationality-framework derivations and boundaries}
\label{sec:p11-ri-proofs}

\begin{annexitem}{annex:p11.ri.individual-joint}{thm:p11-ri-individual-joint}{Individual best responses need not produce a jointly beneficial result}
\begin{proof}
If player 2 chooses \(\mathrm{Coop}\), player 1 gets 3 from
\(\mathrm{Coop}\) and 4 from \(\mathrm{Def}\), so the latter
is strictly better. If player 2 chooses \(\mathrm{Def}\),
player 1 gets 0 and 1 respectively, again strictly favoring
\(\mathrm{Def}\). The same two comparisons hold for player 2
by the symmetric payoff table. Hence \(\mathrm{Def}\) is each
player's unique best response to either action of the other.

Any profile with one player choosing \(\mathrm{Coop}\) permits
that player a strictly beneficial unilateral switch, and thus is
not a mutual-best-response profile. At
\((\mathrm{Def},\mathrm{Def})\), each is already playing the
unique best response, proving existence and uniqueness under the
stated criterion. Its payoffs are \((1,1)\).
The table gives \((3,3)\) at joint cooperation, strictly increasing
both coordinates. Thus the individual criterion does not entail
absence of a joint improvement.

This is an exact consequence of the stipulated table and comparison
criteria. It makes no empirical prediction about what participants
choose and no claim that self-payoff maximization is every
philosophical account's rationality standard.
\end{proof}
\end{annexitem}

\begin{annexitem}{annex:p11.ri.threshold-closure}{thm:p11-ri-threshold-closure}{A fixed threshold below certainty does not ensure conjunctive closure}
\begin{proof}
Since \(1-t>0\), choose an integer
\(n>1/(1-t)\) with \(n\geq2\). The uniform probability
of each singleton in \(\Omega\) is \(1/n\), so
\[
 \mathbb P(E_i)=1-1/n>t\qquad(1\leq i\leq n).
\]
Every \(E_i\) is therefore accepted under the stipulated rule.
No element belongs to every \(E_i\): the element \(j\) is
absent from \(E_j\). Hence their intersection is empty and has
probability zero. Because \(t\geq0\), zero is not strictly
greater than \(t\); the conjunction is not accepted.

Deductive closure under the classical interpretation would include
the conjunction of any finitely many accepted contents. This
specific threshold construction fails that condition. The result
quantifies over every fixed threshold below one, allowing the
finite space to depend on it. It does not claim that every threshold
rule fails on every particular space. A strict threshold of one,
excluded in the statement, would accept nothing; a non-strict rule
at certainty has a different boundary. No psychological frequency
or general impossibility of relating belief and credence follows.
\end{proof}
\end{annexitem}

\begin{annexitem}{annex:p11.ri.no-bridge}{thm:p11-ri-no-bridge}{Independent assertions require a bridge premise}
\begin{proof}
Satisfiability gives a valuation \(v\) of the descriptive atoms
that satisfies every member of \(\Gamma\). Extend it first by
assigning \(N\) true and second by assigning \(N\) false,
leaving all descriptive values fixed. Because \(N\) occurs in
none of the formulas, each extension still satisfies every member
of \(\Gamma\). One extension falsifies \(\neg N\), the other
falsifies \(N\). Thus neither assertion is a consequence of
\(\Gamma\).

For the converse claim, fix \(N\) true and vary \(D\) between
true and false. Both valuations satisfy the premise \(N\),
while respectively falsifying \(\neg D\) and \(D\).
Fixing \(N\) false gives the same two countervaluations for
the premise \(\neg N\). This proves all four failures of
entailment asserted in the statement.

The argument uses precisely the independence supplied by the
valuation space. If an additional bridge axiom restricts that
space, these countervaluations may cease to be available. If
\(\Gamma\) is inconsistent, there is no satisfying valuation
and classical semantic entailment is vacuous, so the first claim
would fail; its satisfiability hypothesis is essential. This is a
logical illustration of missing premises, not a solution to the
metaphysics or semantics of normativity.
\end{proof}
\end{annexitem}

\begin{annexitem}{annex:p11.ri.ellsberg-contradiction}{thm:p11-ri-ellsberg-contradiction}{The common-event contradiction in the introduction's variant}
\begin{proof}
Put \(\Delta u=u(H)-u(L)>0\). The difference between
the expected utility of the \(R\) bet and that of the \(W\)
bet is
\[
 \Delta u\bigl(\mathbb P(R)-\mathbb P(W)\bigr).
\]
Representing the first strict preference therefore requires
\(\mathbb P(R)>\mathbb P(W)\).

The color events are pairwise disjoint. Additivity gives
\[
\begin{aligned}
 \mathbb P(B\cup W)&=\mathbb P(B)+\mathbb P(W),\\
 \mathbb P(B\cup R)&=\mathbb P(B)+\mathbb P(R).
\end{aligned}
\]
Subtracting the expected utility of the second bet from that of
the first gives
\(\Delta u(\mathbb P(W)-\mathbb P(R))\). Thus representing
the second strict preference requires
\(\mathbb P(W)>\mathbb P(R)\). The two strict inequalities
cannot both hold for real numbers. No particular prior over the
unknown color split, nor a particular curvature of utility in
money, was needed.

The contradiction depends on one prior being used for all four
bets, additivity, common state-independent consequence utility,
and representation of strict preference by strictly larger
expected utility. Menu-dependent priors, state-dependent
utilities, nonadditive evaluations or stochastic approximate
fits are not ruled out by this proof. The known red proportion
specifies the source experiment's information but is not itself
the cause of the contradiction. The proof does not establish
how often the pattern is observed or whether its assumptions
should be accepted as a universal norm.
\end{proof}
\end{annexitem}

\begin{annexitem}{annex:p11.ri.majority-cycle}{thm:p11-ri-majority-cycle}{Individual transitivity does not ensure majority transitivity}
\begin{proof}
Each row lists all three alternatives in a strict linear order.
Including the first-to-third comparison makes each row transitive:
voter 1 also ranks \(a\) above \(c\), voter 2 ranks \(b\)
above \(a\), and voter 3 ranks \(c\) above \(b\).

Voters 1 and 3 rank \(a\) above \(b\), giving
\(a\succ_M b\). Voters 1 and 2 rank \(b\) above \(c\),
giving \(b\succ_M c\). Voters 2 and 3 rank \(c\) above
\(a\), giving \(c\succ_M a\). Only voter 1 ranks \(a\)
above \(c\), so \(a\succ_M c\) does not hold. Transitivity
applied to the first two majority comparisons would require it;
the requirement therefore fails.

There are no ties and no individually cyclic input in this
construction. The failure is an exact property of the specified
pairwise-majority rule on this profile. It supplies neither
Arrow's universal impossibility theorem nor an empirical claim
about any electorate, and does not assign cardinal welfare
values to the voters' alternatives.
\end{proof}
\end{annexitem}

\begin{annexitem}{annex:p11.ri.card-check}{thm:p11-ri-card-check}{Which hidden faces can violate the stated rule}
\begin{proof}
Material implication \(\mathsf V(k)\to\mathsf E(k)\) is false
exactly when \(\mathsf V(k)\) is true and \(\mathsf E(k)\)
false. Thus the only violating pattern is a vowel paired with
an odd number.

The visible \(A\) already makes the vowel condition true;
its hidden number can be odd, so its other face must be checked.
The visible \(7\) already makes the even-number condition false;
its hidden letter can be a vowel, so that other face must be checked.
The \(D\) has a false antecedent regardless of its hidden number,
and the \(4\) a true consequent regardless of its hidden letter.
Neither can falsify the stipulated implication.

After inspecting both \(A\) and \(7\), every possible violation
among the four displayed cards is either found or excluded, proving
sufficiency. To prove necessity, if the \(A\) is omitted, choose
an odd hidden number for it, an otherwise nonviolating hidden
letter for \(7\), and any compatible backs for the other cards.
All inspected faces can agree with the rule while the unchecked
\(A\) violates it. If \(7\) is omitted, give it a vowel back
and \(A\) an even back; the unchecked \(7\) then supplies
the corresponding possible violation. These are legitimate hidden
assignments under the two-side setup. Omitting either therefore
prevents a guaranteed check of all displayed possible countercases.

The proof neither certifies any unobserved card nor identifies
an actual person's interpretation or process. Different task
instructions, deontic meanings, information goals or sampling
assumptions require an appropriate new criterion.
\end{proof}
\end{annexitem}
